\documentclass[ecta,nameyear,final,supplement]{econsocart}
\RequirePackage[colorlinks,citecolor=blue,linkcolor=blue,urlcolor=blue]{hyperref}
\usepackage{amssymb,bm,booktabs,array,multirow,graphicx}

\startlocaldefs
\allowdisplaybreaks
\renewcommand{\topfraction}{0.9}\renewcommand{\bottomfraction}{0.7}\renewcommand{\textfraction}{0.08}\renewcommand{\floatpagefraction}{0.75}

\newcommand{\E}{{\mathbb E}}
\newcommand{\PP}{{\mathbb P}}
\newcommand{\dr}[1]{{\mathrm d}#1}
\newcommand{\iu}{{\mathrm i}}
\newcommand{\Rea}{\operatorname{Re}}
\newcommand{\Ima}{\operatorname{Im}}
\newcommand{\Var}{\operatorname{Var}}
\newcommand{\Cov}{\operatorname{Cov}}
\newcommand{\erf}{\operatorname{erf}}
\newcommand{\erfc}{\operatorname{erfc}}
\newcommand{\tr}{\operatorname{tr}}
\newcommand{\adj}{\operatorname{adj}}
\newcommand{\sech}{\operatorname{sech}}
\newcommand{\diag}{\operatorname{diag}}
\newcommand{\phih}{\hat\phi}
\newcommand{\tauh}{\hat\tau}
\newcommand{\Fk}{\textsf{\upshape DK}}
\newcommand{\Fu}{\textsf{\upshape DU}}
\newcommand{\Ft}{\textsf{\upshape DT}}
\newcommand{\Sk}{\textsf{\upshape SK}}
\newcommand{\Su}{\textsf{\upshape SU}}
\newcommand{\St}{\textsf{\upshape ST}}
\newcommand{\cF}{\mathcal{F}}
\newcommand{\cP}{\mathcal{P}}
\newcommand{\cQ}{\mathcal{Q}}
\newcommand{\cG}{\mathcal{G}}
\newcommand{\bX}{\boldsymbol{X}}
\newcommand{\bQ}{\boldsymbol{Q}}
\newcommand{\bK}{\boldsymbol{K}}
\newcommand{\bI}{\boldsymbol{I}}
\newcommand{\bGam}{\boldsymbol{\Gamma}}
\newcommand{\bSig}{\boldsymbol{\Sigma}}
\newcommand{\bP}{\boldsymbol{P}}
\newcommand{\bG}{\boldsymbol{G}}

\theoremstyle{plain}
\newtheorem{theorem}{Theorem}
\renewcommand{\thetheorem}{S\arabic{theorem}}
\theoremstyle{definition}
\newtheorem{remark}{Remark}

\endlocaldefs

\begin{document}

\begin{frontmatter}
\title{Supplement to ``The Exact Limiting Distribution of the Dickey--Fuller $t$-Ratio'': the conditional integral representation and the numerical evaluation}
\runtitle{Supplement: the conditional integral and the numerical evaluation}

\begin{aug}
\author[purdue]{\fnms{Yong}~\snm{Bao}\ead[label=e1]{ybao@purdue.edu}}
\author[hslu]{\fnms{Simon A.}~\snm{Broda}\ead[label=e2]{simon.broda@hslu.ch}}
\author[rotman,hkust]{\fnms{Raymond}~\snm{Kan}\ead[label=e3]{Raymond.Kan@rotman.utoronto.ca}}
\address[purdue]{\orgdiv{Department of Economics}, \orgname{Purdue University}}
\address[hslu]{\orgdiv{Institute of Financial Services IFZ}, \orgname{Lucerne University of Applied Sciences and Arts}}
\address[rotman]{\orgdiv{Rotman School of Management}, \orgname{University of Toronto}}
\address[hkust]{\orgdiv{School of Business and Management}, \orgname{Hong Kong University of Science and Technology}}
\end{aug}

\begin{abstract}
This supplement gives a second representation of the limiting density and distribution function of the Dickey--Fuller $t$-ratio, the same double integral as in the paper taken in the other order, with the inner inversion along saddle-point paths (Section~\ref{sec:double}, with the proof of Theorem~\ref{th:double}); the numerical evaluation of both representations, their cost and accuracy, and the programs (Section~\ref{sec:numerics}); and the checks against the closed forms, the moments, the critical values and the simulations in the literature.  Numbers without the prefix S refer to the paper.
\end{abstract}
\end{frontmatter}

\section{The conditional integral representation: the other order of integration}
\label{sec:double}

The representation of Theorem~2 of the paper, which we call the line integral representation here, and the conditional integral representation of this section are one double integral, taken in its two orders: this section conditions on $V$.  Numbers without the prefix S refer to the paper.  Let $R=U+\frac12$ and $a(v)=\frac12+x\sqrt v+c_0v$, so that $\{\tau\le x\}=\{R\le a(V)\}$.  By (22), under the tilt $e^{tV}$ the measure of $R$ is $P(t)\varphi_t$, so that $P(t)\Psi_t(a)=\E[e^{tV}\mathbf{1}\{R\le a\}]$ is the Laplace transform in $v$ of the measure $\E[\delta(V-v)\mathbf{1}\{R\le a\}]$, and for real $\sigma<t_*$
\begin{equation}
\begin{aligned}
\cF(x)&=\PP\{R\le a(V)\}=\int_0^\infty H(v;x)\dr{v},\\
H(v;x)&=\E\big[\delta(V-v)\,\mathbf{1}\{R\le a(v)\}\big]=\frac1{2\pi\iu}\int_{\sigma-\iu\infty}^{\sigma+\iu\infty}P(t)\,\Psi_t\big(a(v)\big)\,e^{-tv}\dr{t}.
\end{aligned}
\label{eq:limrep}
\end{equation}
With the $v$-integral taken inside, on one line $\Rea t=\tau$ for every $v$, the substitution $v=w^2$ gives $\int_0^\infty\Psi_t(a(v))e^{-tv}\dr{v}=\Phi_x(t)/t$, and (\ref{eq:limrep}) is the distribution function of Theorem~2; the density is the same with $\sqrt v\,\varphi_t(a(v))$ in place of $\Psi_t(a(v))$.  This section takes the $t$-integral inside.  Nothing is gained or lost in dimension: $\Phi_x(t)$ and $H(v;x)$ are the two partial integrals of one function of $(t,v)$, and both are evaluated by quadrature.  What changes is the contour.  Inside the $v$-integral the line in $t$ is fixed, and the integrand of $\Phi_x$ along $v$ is a Gaussian transform of a density that is explicit, or with a fixed start a one-dimensional integral, whose analytic structure is known once and for all.  Inside the $t$-integral the line may be moved, for each $v$, to a path through the saddle point of $\log M_V(t)-tv$, on which the integrand decays fast whatever $x$, at the price of a new path for every $v$.  This is the route of the finite-sample theory in \cite{BaoBrodaKan2026}, where it is the natural one, and in the limit it is an independent computation of everything: it shares with Theorem~2 only Proposition~1, in different forms, and nothing of the quadrature.
Under the tilt $e^{tV}$, for real $t<t_*$ and by analytic continuation on the contours, $R$ is $\kappa_1Z_1^2+\kappa_2Z_2^2$ with the weights $\kappa_k=1/(2s_k)$ of Lemma~1, and its normalized density is $(-s_1s_2)^{1/2}\varphi_t$, respectively $s_1^{1/2}\varphi_t$; so $M_V(t)$ times that density is $P(t)\varphi_t(r)$, and the joint density of $(R,V)$ is
\begin{equation}
h(r,v)=\frac1{2\pi\iu}\int_{\sigma-\iu\infty}^{\sigma+\iu\infty}P(t)\,\varphi_t(r)\,e^{-tv}\dr{t},
\label{eq:hrv}
\end{equation}
for any real $\sigma<t_*$.

\begin{theorem}[The conditional integral representation]
\label{th:double}
Let $\sigma<t_*$ be real.  (i) For every $v>0$,
\begin{equation}
H(v;x)=\frac1{2\pi\iu}\int_{\sigma-\iu\infty}^{\sigma+\iu\infty}M_V(t)\,e^{-tv}\,\cP\big(\kappa_1(t),\kappa_2(t);a(v)\big)\dr{t},
\label{eq:Hlim}
\end{equation}
where, with $\phi(l;\omega^2)=(2\pi\omega^2)^{-1/2}e^{-l^2/(2\omega^2)}$ and $l_0=\sqrt{\max(0,-2a)}$,
\begin{equation}
\cP(\kappa_1,\kappa_2;a)=2\int_{l_0}^\infty\phi(l;-2\kappa_2)\,\erf\Big(\sqrt{\frac{2a+l^2}{4\kappa_1}}\Big)\dr{l}
\label{eq:Pkappa}
\end{equation}
in cases \Fu\ and \Ft, and $\cP(\kappa_1,0;a)=\erf(\sqrt{a/(2\kappa_1)})$ for $a>0$ and $0$ for $a\le0$ in case \Fk; (\ref{eq:Pkappa}) requires $\Rea\kappa_1>0>\Rea\kappa_2$ on the line of integration, which Lemma~1 provides.  The line may be replaced by any path $\Gamma$ that crosses the real axis once, at $\sigma$, lies in $\Rea t\ge\sigma$ and tends to infinity in the lower and the upper half-plane at its two ends, provided the sign conditions hold on $\Gamma$ and on the region between $\Gamma$ and the line, and the integrals of the integrand over the parts of the circles $|t|=\varrho$ between $\Gamma$ and the line tend to zero as $\varrho\to\infty$.

(ii) With $h$ as in (\ref{eq:hrv}) and $a'(v)=x/(2\sqrt v)+c_0$, the density and distribution function of the limit are the double integrals
\begin{align}
f(x)&=\int_0^\infty\sqrt v\;h\big(a(v),v\big)\dr{v},
\label{eq:flim}\\
\cF(x)&=\cP\big(\kappa_1(0),\kappa_2(0);\tfrac12\big)+\int_0^\infty a'(v)\,G\big(a(v),v\big)\dr{v},
\label{eq:Flim}\\
G(r,v)&=\frac1{2\pi\iu}\int_{\sigma'-\iu\infty}^{\sigma'+\iu\infty}P(t)\,\varphi_t(r)\,\frac{e^{-tv}}t\dr{t},
\label{eq:Glim}
\end{align}
where $0<\sigma'<t_*$; the first term is $\PP(R\le\frac12)$, and $G(r,v)=\int_v^\infty h(r,v')\dr{v'}$.  The lines may be replaced by paths $\Gamma$ as in (i); for $G$, a path that crosses the real axis at a negative point adds the residue $(-s_1s_2)^{1/2}\varphi_0(r)$, the density of $R$ without tilt.
\end{theorem}

For real weights $\kappa_1>0>\kappa_2$, $\cP$ is the probability that $\kappa_1Z_1^2+\kappa_2Z_2^2\le a$, conditioning on $Z_2$; for the complex weights on the line of integration it is the analytic continuation of that probability in $(\kappa_1,\kappa_2)$, with no branch ambiguity, since the square roots in (\ref{eq:Pkappa}) are of nonnegative reals and of numbers with positive real part.  Part (ii) replaces the probability $\cP$, itself an integral, by the density of $R$ under the tilt, so that both the density and the distribution function are double integrals with closed-form integrands.  The density (\ref{eq:flim}) integrates the joint density of $(R,V)$ along the curve $r=a(v)$, as in the construction of \cite{DanielsYoung1991} for the Studentized mean; the distribution function (\ref{eq:Flim}) follows from (\ref{eq:limrep}) by an integration by parts in $v$.  Integrating over $v$ first instead, as in (\ref{eq:limrep}), gives back Theorem~2, with $\Phi_x$ and $\Phi^f_x$ as the inner integrals.  In case \Fk\ the probability $\cP$ is an error function, so part (i) is already a double integral; we use it there, because at the support edge $a(v)=0$ the density $\varphi_t$ has an inverse square-root singularity and part (ii) is then accurate only to about $10^{-8}$.  In cases \Fu\ and \Ft\ we use part (ii).  The proof is at the end of this section.

\paragraph{The initial observation}  With a fixed start, $R$ is the noncentral form of Proposition~2 under the tilt, and the theorem holds with $M_{\gamma_0}(0,t)$ in place of $M_V$, the weights unchanged, and $\cP$ replaced by
\begin{equation}
\cP(\kappa_1,\kappa_2;a)=2\int_{l_0}^\infty\bar\phi(l;m_2,-2\kappa_2)\,\bar F\big((2a+l^2)^{1/2};m_1,2\kappa_1\big)\dr{l},
\label{eq:Pnc}
\end{equation}
where
\[
\begin{gathered}
\bar\phi(l;m,w)=\frac12\{\phi(l-m;w)+\phi(l+m;w)\},\\
\bar F(r;m,w)=\frac12\erf\{(r-m)/(2w)^{1/2}\}+\frac12\erf\{(r+m)/(2w)^{1/2}\},
\end{gathered}
\]
the probability that $X^2\le r^2$ for $X\sim N(m,w)$; in case \Fk, $\cP=\bar F(\tilde a;m,w)$ with $\tilde a=(2a+\gamma_0^2)^{1/2}$, $m=\gamma_0/F_0$ and $w=F_1/F_0$, for $2a+\gamma_0^2>0$, and $\cP=0$ otherwise.  Both are even in the means, so only $m_k^2=\gamma_0^2\mu_k$ enters, and they are analytic in the weights and the squared means under the sign conditions.  Part (ii) holds with the density of $R$ under the tilt in place of $(-s_1s_2)^{1/2}\varphi_t$: in case \Fk\ with a fixed start it is $\{\phi(\tilde r-m;w)+\phi(\tilde r+m;w)\}/\tilde r$ with $\tilde r=(2r+\gamma_0^2)^{1/2}$ and the same $m$ and $w$, on $2r+\gamma_0^2>0$, and with a stationary start it is (23) with the roots of $D_{\rm S}$.  So five of the six configurations are double integrals with closed-form integrands: case \Fk\ with a zero or a fixed start through part (i), and the stationary starts and cases \Fu\ and \Ft\ with a zero start through part (ii).  In cases \Fu\ and \Ft\ with a fixed start the density of $R$ under the tilt is (33), itself an integral (the density of a product of two correlated noncentral normal variables, which has no closed form: a double series of Whittaker functions, by expanding each noncentral $\chi^2_1$ as a Poisson mixture of central $\chi^2_{1+2j}$); there part (i) is used with (\ref{eq:Pnc}), a triple integral, with the squared means from (10).  With a fixed start the support edge of case \Fk\ moves to $a(v)=-\gamma_0^2/2$, where $J(1)=0$, and in the other cases the kink of the density of $R$ stays at $a(v)=0$, where $X_1^2=X_2^2$.

\paragraph{Analyticity beyond $t_*$}  The integrands of the conditional integral with a zero start are analytic beyond $t_*$ by the Girsanov forms of Section~3.1.  In case \Fk, for $a>0$,
\begin{equation}
M_V(t)\,\cP(\kappa_1,0;a)\,e^{-tv}=e^{-c/2}\Big(\frac{\vartheta}{\sinh\vartheta}\Big)^{1/2}\frac{\erf\{(a\eta)^{1/2}\}}{\eta^{1/2}}\,e^{-tv},
\label{eq:dkprod}
\end{equation}
and the factor $\eta^{-1/2}$ of $M_V$, which is infinite at $t_*$, where $\eta=0$, cancels: $\erf(w)/w$ is an entire function of $w^2$, so (\ref{eq:dkprod}) is analytic up to the first zero of $\sinh\vartheta/\vartheta$, at $t=(c^2+\pi^2)/2$.  In cases \Fu\ and \Ft\ the product $P(t)\varphi_t(r)$ of (15) is analytic up to $(c^2+\pi^2/4)/2$ at most, where $\cosh\vartheta=0$.  This matters for $c>0$.  There the saddle point of the whole integrand on the real axis often lies beyond $t_*$, while a path through the saddle point of $\log M_V(t)-tv$ alone must cross the axis to the left of $t_*$; along such a path the factor $\cP$ or $\varphi_t(a(v))$ varies over many orders of magnitude, and the integral loses accuracy (in case \Fk\ with $c_0=0$ the error of that route at $x=0$ is $2\times10^{-9}$ at $c=2$ and $5\times10^{-3}$ at $c=5$, where the probability is $0.0170$; in cases \Fu\ and \Ft\ it fails at $c=5$).  With a zero start we therefore integrate, for each $v$, along a path through the saddle point of the logarithm of the whole integrand.  With an initial observation the integrands are not of this form, and we integrate along a parabola through the saddle point of $\log M_V(t)-tv$, with the sign conditions checked at every node; this route is accurate to $10^{-11}$ for $c\le2$, and the programs warn when their error estimates exceed $10^{-10}$, as they do by $c=3$ with a fixed start in cases \Fu\ and \Ft.  The line integral of Theorem~2 has no such limitation, with a fixed start as well, because its contour is the fixed line $\Rea t=\tau<t_*$ on which the factorization of Lemma~1 holds; it is the route we take, and the conditional integral takes over far in the tails, where the error of the line integral, absolute rather than relative, shows, and at a nonzero local null value, where it is the faster one (Section~\ref{sec:numerics}).

\paragraph{Proof of Theorem~\ref{th:double}}
By (4), $R$ is a quadratic form $R=\frac12\bm\ell'\bm Q\bm\ell$ in a vector $\bm\ell$ of at most three Gaussian functionals of $W$, of rank one or two.  Approximate $W$ on a grid of $N$ points: $V$ becomes a quadratic form $\bm x'\bm K_V\bm x$ in $\bm x\sim N(\bm0,\bI_N)$, with $\bm K_V$ positive semidefinite, and $\bm\ell=\bm\Gamma\bm x$.  Let $t=\sigma+\iu y$ with $\sigma$ real and $\bP=\bI_N-2t\bm K_V$, whose real part $\bI_N-2\sigma\bm K_V$ is positive definite when $\sigma<t_*$.  For bounded measurable $h$, completing the square in the Gaussian integral gives
\[
\E\big\{e^{tV}h(\bm\ell)\big\}=|\bP|^{-1/2}\int h(\bm l)\,\phi_{\bSig}(\bm l)\dr{\bm l},\qquad\bSig=\bm\Gamma\bP^{-1}\bm\Gamma',
\]
where $\phi_{\bSig}$ is the Gaussian density formula evaluated at the complex covariance $\bSig$, a complex-valued function.  It is integrable because $\Rea\bSig^{-1}$, a Schur complement of the accretive matrix $\bP$ after a change of coordinates, is positive definite.  With $h=\mathbf{1}\{R\le a\}$ the integral is an analytic function of $\bSig$ on this set.  For real $\bSig$ it is the probability that $\kappa_1Z_1^2+\kappa_2Z_2^2\le a$, where $\kappa_1$, $\kappa_2$ are the nonzero eigenvalues of $\frac12\bSig\bm Q$, which equals (\ref{eq:Pkappa}), as the integral over $Z_1$ given $Z_2$ shows.  The right side of (\ref{eq:Pkappa}) is analytic in $(\kappa_1,\kappa_2)$ where $\Rea\kappa_1>0>\Rea\kappa_2$, so the two agree there by analytic continuation.  The moment generating function of $R$ under the tilt, $\prod_{k=1}^2(1-2s\kappa_k)^{-1/2}$, identifies the $\kappa_k$ with $1/(2s_k)$, by Lemma~1.  Letting $N\to\infty$, the finite-dimensional quantities converge to those of the limit, locally uniformly in $t$, because the moment generating functions do.  Finally $\E\{e^{tV}\mathbf{1}(R\le a)\}$ is the Laplace transform of the measure $\E\{\,\cdot\,;R\le a\}$ in $v$, and (\ref{eq:Hlim}) is its inversion on the line $\Rea t=\sigma$.  For the path $\Gamma$, the integrand of (\ref{eq:Hlim}) is analytic on the region between $\Gamma$ and the line: the singularities of $M_V$ are the zeros of $D(0,\cdot)$, the Fredholm determinant of the positive semidefinite covariance kernel of $V$, which are real; the weights are analytic wherever they are distinct, which the sign conditions ensure, and $\cP$ is analytic in them there.  By assumption, the integrals over the arcs of the circles $|t|=\varrho$ that join $\Gamma$ and the line tend to zero as $\varrho\to\infty$; Cauchy's theorem on the bounded regions cut off by these arcs then moves the line to $\Gamma$.

For part (ii), take $h$ in the Gaussian integral to be a function of $R$ alone.  For real $\bSig$ the law of $R$ is that of $\kappa_1Z_1^2+\kappa_2Z_2^2$, whose density is $(-s_1s_2)^{1/2}\varphi_t$ (respectively $s_1^{1/2}\varphi_t$): in case \Fk\ the scaled $\chi^2_1$ density, and in cases \Fu\ and \Ft\ the convolution of the densities of $\kappa_1Z_1^2$ and $\kappa_2Z_2^2$, which reduces with $K_0(w)=\int_0^\infty e^{-w\cosh q}\dr{q}$ to (23).  Both sides are analytic in $(\kappa_1,\kappa_2)$ where $\Rea\kappa_1>0>\Rea\kappa_2$, so, as for part (i), $\E\{e^{tV}h(R)\}=P(t)\int_{-\infty}^{\infty}h(r)\varphi_t(r)\dr{r}$ on the line, and in the limit $N\to\infty$.  Hence $P(t)\varphi_t(r)$ is the Laplace transform in $v$ of the joint density of $(R,V)$, and (\ref{eq:hrv}) is its inversion.  The integral converges absolutely for $r\ne0$: with $\Rea s_1>0>\Rea s_2$, $|e^{-(s_1+s_2)r/2}K_0((s_1-s_2)|r|/2)|$ behaves like $|e^{-s_1r}|$ for $r>0$ and like $|e^{-s_2r}|$ for $r<0$ when $|t|$ is large, so $\varphi_t(r)$ grows at most like a power of $|t|$, while $P$ decays like $\exp(-\frac12|\Ima t|^{1/2})$.  Then $H(v;x)=\int_{-\infty}^{a(v)}h(r,v)\dr{r}$, and differentiating $\int_0^\infty H(v;x)\dr{v}$ in $x$, with $\partial a/\partial x=\sqrt v$, gives (\ref{eq:flim}).  For the distribution function, $G(r,v)=\int_v^\infty h(r,v')\dr{v'}$ is the density of $R$ on the event $\{V>v\}$, and for $\sigma'>0$ the integral of (\ref{eq:hrv}) over $v'>v$ may be taken under the $t$-integral, which gives (\ref{eq:Glim}).  Moving its line across the simple pole at $t=0$, where $M_V(0)=1$, gives the residue $(-s_1s_2)^{1/2}\varphi_0(r)$.  With $\Psi(v)=\int_{-\infty}^{a(v)}G(r,v)\dr{r}=\PP\{R\le a(v),\,V>v\}$ and $\partial G/\partial v=-h$,
\[
\Psi'(v)=a'(v)\,G\big(a(v),v\big)-H(v;x).
\]
Since $\Psi(0)=\PP(R\le\frac12)$, $\Psi(v)\to0$ as $v\to\infty$, and $G(r,v)$ is bounded near $v=0$ for $r=a(v)$ near $\frac12$, so that $a'(v)G(a(v),v)$ is integrable at $v=0$, integrating over $v\in(0,\infty)$ gives (\ref{eq:Flim}).  At $t=0$ the weights are real, and $\PP(R\le\frac12)=\cP(\kappa_1(0),\kappa_2(0);\frac12)$ by part (i).  With a fixed start the Gaussian integral has the deterministic part of $\bm\ell$ carried along, $\E\{e^{tV}h(\bm\ell)\}=|\bP|^{-1/2}e^{\gamma_0^2Q(0,t)}\int h(\bm l)\phi_{\bSig}(\bm l-\bm m)\dr{\bm l}$ with a complex mean $\bm m$ proportional to $\gamma_0$, and for real $\bSig$ and $\bm m$ the integral with $h=\mathbf{1}\{R\le a\}$ is $\PP\{\frac12(X_1^2-X_2^2)\le a\}$ for independent $X_k\sim N(m_k,\pm2\kappa_k)$, which is (\ref{eq:Pnc}) after conditioning on $X_2$ and replacing $X_2$ by $|X_2|$; the right side of (\ref{eq:Pnc}) depends on the means only through $m_k^2$ and is analytic in the weights and in $m_k^2$, so analytic continuation and the limit $N\to\infty$ give the result as before, with $m_k^2=\gamma_0^2\mu_k$ by Proposition~2.  With a stationary start $R$ is a central form of rank two in the Gaussian vector that includes $\gamma_0$, and the theorem applies with $D_{\rm S}$.  \qed


\section{Numerical evaluation}
\label{sec:numerics}

Both representations are evaluated to eleven or twelve digits by the MATLAB programs \texttt{ltratio\_cdf}, \texttt{ltratio\_pdf}, \texttt{ltratio\_icdf} and \texttt{ltratio\_mom} that accompany the paper, and by the implementations in Julia and in C, the latter with the same MATLAB interface.  Either representation can be selected by an option; by default the programs choose between them point by point, by their own timings.  At $c_0=0$ they take the line integral.  At a nonzero local null value the conditional integral is the faster one, by a factor of two to ten, except for $c_0>0$ at $x\ge0$, where the line integral has no breakpoint; in cases \Su\ and \St\ the MATLAB programs keep the line integral at $x<0$, because their conditional integral with a stationary start costs half a second a value there.  Far in the tails the conditional integral takes over in any case, because the error of the line integral is absolute rather than relative (the paragraph on accuracy below); conversely, where the conditional integral reports that its error estimate is not met, the programs return the value of the line integral.  With a fixed nonzero start the MATLAB programs use the conditional integral throughout, since it is their only route; the Julia and C implementations carry the noncentral kernel of the line integral and use it there.  This section describes the two computations.

\subsection{The line integral representation}
\label{app:single}

\paragraph{The line integral representation}  The line of (25) is $\Rea t=\tau$ with $\tau=1$ for $c\le0$ and a zero start, where $t_*$ exceeds $1$ in the three cases ($1.23$, $4.93$ and $19.7$ at $c=0$, increasing as $c$ decreases), and $\tau=\min(1,t_*/2)$ otherwise ($t_*$ is $0.046$, $0.10$ and $0.97$ at $c=3$, $4\times10^{-7}$ to $10^{-6}$ at $c=10$, and $15.1$, $23.8$ and $37.8$ with a stationary start at $c=-5$).  The integrand is conjugate symmetric on the line, so the integral is twice the real part of the integral over $\Ima t=w>0$.  Its two parts are treated differently.  The part of $P(t)\Phi_x(t)/t$ without the oscillating factors is analytic in a strip around the line and decays like $\exp(-\frac12\sqrt w)$, the decay of $P$: with $w=\tau\sinh u$ the trapezoidal rule in $u$ converges geometrically, and $123$ nodes with the step $0.075$, up to $w=4900$, give $10^{-15}$.  Each part with a factor $e^{-tv_j^2}$, $v_j$ a root of $X(v)$ ($v_j=1/(2|x|)$ at $c_0=0$ and $x<0$), is written as $G_j(w)e^{-\iu wv_j^2}$ with $G_j$ smooth and integrated by Filon's method: $G_j$ is interpolated by Legendre polynomials on panels that grow geometrically from the scale $\min(\frac12,1/v_j^2)$ of the pole at $t=0$, with $20$ Gauss--Legendre nodes each, and the oscillation is integrated exactly, $\int_{-1}^1P_n(x)e^{-\iu\Omega x}\dr{x}=2(-\iu)^nj_n(\Omega)$.  The parts are separated inside the transforms (26) by an entire partition of unity, so that each remains analytic in $t$, and the transforms themselves are evaluated by Gauss--Legendre rules on paths in the complex $v$-plane on which the Gaussian factor decays instead of oscillating, with the function $e^zK_0(z)$ evaluated by series and asymptotic expansions to $10^{-14}$.  For $c>0$, and with $c_0<0$ in the upper tail, the exponent $-tv^2-s_dX(v)$ of the integrand has a saddle point whose phase grows like $x^2$ along the line; where its change between the nodes of the trapezoidal rule exceeds $\frac14$, the smooth part is integrated instead on geometric panels with Gauss--Legendre rules whose node counts follow the phase.  With $c_0\ne0$ the roots of $X(v)$, and for $c_0>0$ the vertex of the parabola, are breakpoints of the inner integral, and the paths between them are straight pieces along the directions of steepest descent of the exponent at their ends, joined through the real axis where a direct connection would cross a region in which the Gaussian factor grows.


\paragraph{The transforms}  The transforms (26) are integrals of $\varphi_t$ against Gaussian kernels over half-lines and finite segments of the $v$-axis.  At $c_0=0$ in case \Fk\ the segment integrals are written as differences of half-line integrals, which are the special functions (27); otherwise the Gaussian transforms are evaluated by quadrature as follows, and the same rules, with the $K_0$ kernel replaced by $X^{-1/2}$, evaluate (27).  For $x<0$ and $c_0=0$ the integrals over $(0,\frac12)$ in $X$ contain the contributions of the support edge $X=0$, which carry the factor $e^{-ty}$, $y=1/(4x^2)$, and oscillate in $\Ima t$ with the period $8\pi x^2$, and the contributions of the neighborhood of $X=\frac12$, where the Gaussian factor is centered, which are analytic in $t$ without oscillation.  They are separated by a partition of unity, $\chi(X)=\frac12[1+\erf\{\kappa(\frac14-X)\}]$ with $\kappa=24$, applied to the integrand over $(0,\frac12)$: $\chi$ is entire, $1$ and $0$ at the ends of the segment to $10^{-16}$, so that both parts are integrals of analytic functions between the same real endpoints, hence analytic and conjugate symmetric in $t$, and only the part with $\chi$ carries the factor $e^{-ty}$.  The partition is applied to the complex path point on a fixed path, so that its definition does not depend on $t$.  In case \Fk\ the same partition, in the variable $v$ with $\kappa=12/v^*$, $v^*=1/(2|x|)$, reproduces the split of (29) up to terms of relative size $10^{-16}$.

\paragraph{The paths}  The half-lines from $X=\frac12$ and from the edge $X=0$ outward are rotated by $-\frac12\arg t$, on which the Gaussian factor decays without oscillating, and the segment $(0,\frac12)$ is bulged into the half-plane $\Ima(X-\frac12)^2<0$ along a fixed curve with unit slope at both ends, the direction of steepest descent there, and height $\frac1{16}$, which keeps $|\chi|\le e^{9/4}$ on it.  The endpoint singularities, $X^{-1/2}$ in case \Fk\ and $\log X$ otherwise, are removed by the substitution $X=\delta u^4$; each path is divided at the scales of its factors, the width $|x|(2\Ima t)^{-1/2}$ of the Gaussian, the decay lengths of the exponentials and the distance $1/\kappa$ over which $\chi$ varies, with $32$ to $96$ nodes per piece and the pieces growing geometrically along the half-lines.  The constant part of the kernel in the cases with the $K_0$ kernel is the total mass $(-s_1s_2)^{-1/2}$ minus a short integral over $X>\frac12$ on the ray $X=\frac12+re^{-\iu\arg s_1}$, on which $\varphi_t$ decays without oscillating, or, when $|s_1|$ is small, the integral over $X<\frac12$ with the ray $X=-re^{-\iu\arg(-s_2)}$ from the edge.  The branches of $P(t)$, $s_1(t)$, $s_2(t)$ and $(-s_1s_2)^{-1/2}$ are fixed by continuity from the real point $t=\tau$ along each set of nodes; for $c>2$ the roots and $P(t)$ come from the Girsanov forms (15), and the smaller root in modulus from the product of the roots rather than the quadratic formula, which loses it to cancellation when the roots differ by orders of magnitude (at $c=10$ they are about $-100$ and $3\times10^{-8}$).

\paragraph{The Bessel function}  The function $e^zK_0(z)$, needed at points with $\Rea z\ge0$ only, is evaluated to $10^{-14}$ by its power series for $|z|<2$, by Taylor expansions about the centers of cells of side $\frac12$ for $|z|<25$, the coefficients following from $K_0$ and $K_1$ at the center through Bessel's equation, and by its asymptotic series beyond, at a tenth of the cost of a library routine.

\paragraph{The outer integral}  The part without oscillating factors is integrated by the trapezoidal rule in $u$, $w=\tau\sinh u$, with the step $0.075$ up to $w=4900$ ($123$ nodes).  Each part with a factor $e^{-tv_j^2}$ is integrated by Filon's method on panels that grow geometrically from $\min(\frac12,1/v_j^2,\tau/2)$, the scale of the pole at $t=0$ and of the transient of the transition region of $\chi$, which oscillates at the frequency $3v_j^2/4$ and dies out over $300/v_j^2$; below $w=300/v_j^2$ the panels are at most $25/v_j^2$ long, $G_j$ is interpolated at $20$ Gauss--Legendre nodes by Legendre polynomials, and the moments $\int_{-1}^1P_n(x)e^{-\iu\Omega x}\dr{x}=2(-\iu)^nj_n(\Omega)$ integrate the oscillation exactly; $18$ to $25$ panels reach $w=4900$.  This part is omitted when $3\tau v_j^2/16>40$, where it is below $10^{-17}$.

\paragraph{The saddle factor}  For $c>0$ the roots $s_1$, $s_2$ at $t=\tau$ approach zero as $c$ grows, and the Gaussian parts of the transforms, of the form $\int_0^\infty A(u)e^{-tu^2/x^2-su}\dr{u}$ with $A$ algebraic and $s=s_1$ or $s_2$, grow or oscillate along the ray $u=re^{-\iu\arg(t)/2}$ for $x>0$ large: the model integral with $A=1$ is $\frac12x(\pi/t)^{1/2}e^{z^2}\erfc(z)$ with $z=sx/(2t^{1/2})$, which for $\Rea z<0$ contains the term $2e^{z^2}$ of the saddle point $u_0=-sx^2/(2t)$ of the exponent.  For $\Rea z<0$ these parts are taken along the segment from $0$ to $u_0$ and the line of steepest descent from $u_0$, on which the exponent is $z^2\lambda(2-\lambda)$, $0\le\lambda\le1$, and $z^2-|t|r^2/x^2$, and otherwise along the level curve $\Ima(tu^2/x^2+su)=0$ from $0$, on which the exponent is $-v^2$ for real $v$ and $u=2v^2/\{s+(s^2+4tv^2/x^2)^{1/2}\}$, which ends in the valley of the positive real axis precisely when $\Rea z\ge0$.  The factor $e^{z^2}$, present for $c>0$ beyond a small value of $w$ and never for $c\le0$, where $\Rea z\ge0$ along the line, oscillates in $w$ with a phase that grows like $x^2$; where its change between the nodes of the trapezoidal rule exceeds $\frac14$, the smooth part is integrated instead on the geometric panels with Gauss--Legendre rules whose node counts follow the phase.  The cost then grows with $x^2$ in the upper tail.

\paragraph{The local null value}  With $c_0\ne0$ the breakpoints of the inner integral in $v$ are the positive roots of $X(v)$ and, for $c_0>0$ and $-\sqrt{2c_0}<x<0$, the vertex $v_0=-x/(2c_0)$ of the parabola, where $X$ has the complex roots $v_0\pm\iu\ell$, $\ell=(X(v_0)/c_0)^{1/2}$, from which the cuts of $X^{-1/2}$ or $\log X$ run vertically away from the real axis; the paths leave $v_0$ within $45$ degrees of the real axis, so that they pass the branch points at the distance $0.7\ell$ at least, the first pieces of the quadrature at $v_0$ have the length $\ell$, and the ray from $v_0$ is straight.  The first piece at a root $v_j$ has the length $\ell_j/2$, $\ell_j$ the distance of the nearest other root, and the first piece at $v=0$ the length $\ell_0/2$, $\ell_0$ the distance of the nearest root from $0$ (for $c_0<0$ and large $x$ the negative root $1/(2c_0v_1)\approx-1/(2x)$ lies close to $0$).  On a finite segment between two breakpoints $a<b$ the integrand is $W(v)e^{E(v)}\varphi^{\rm alg}_t(X(v))$ with the exponent $E(v)=-tv^2-s_dX(v)=-\alpha v^2-\beta_0v-\frac12s_d$, $\alpha=t+c_0s_d$, $\beta_0=s_dx$, a quadratic in $v$, and the path is three straight pieces: from $a$ and from $b$ in the directions of steepest descent of $E$ at the ends, each kept within $80$ degrees of the segment and stopped where the exponent has fallen by $60$, where its quadratic part turns it up, at the depth $L/8$, or at the half-length $L/2$; and the connection between their ends, which is the straight line unless the exponent along it rises above its values at the ends, and otherwise the polyline through the real axis, on which the exponent is that of the integrand itself (when $\Rea\alpha\approx0$ and $\Ima\beta_0$ is large the two descent pieces end on opposite sides of the real axis, and a straight connection crosses the quadrant in which $e^{-\alpha v^2}$ grows by $e^{17}$ and more).  Every piece is integrated by Gauss--Legendre rules on sub-pieces graded toward the ends, with the node counts set by the change of $E$ in modulus and phase at $1.2$ nodes per unit plus $16$ for a stationary point of the phase inside, and sub-pieces on which the exponent is below $\max\{E(a),E(b)\}-58$ are skipped; $X$ is formed from the offset from the nearer root, so that it is exact near the roots.  The exponent has its saddle at $v^*=-\beta_0/(2\alpha)$, where it exceeds its value at $0$ by $\beta_0^2/(4\alpha)$.  With $c_0<0$ and large $x$ the saddle sits inside the segment at $v^*\approx b/4$, with a phase of about $0.005x^2\Ima t$ in cases \Fu\ and \Ft\ at $c=5$, which the trapezoidal rule cannot follow; the phase is therefore computed also when a breakpoint is present, whenever $0<\Rea v^*<v_1$, and the smooth part takes the Gauss--Legendre route of the saddle factor, with the partition $\chi$ between the smooth and the edge parts placed midway between $v^*$ and $v_1$ and sharpened to $\kappa=12/(v_1-v^*)$, so that the saddle leaks into the edge part by at most $\erf(-6)$, and the depth of the paths reduced to $1.5/\kappa$.  The saddle-factor part needs about three outer nodes per radian of the phase over the range in which it is not negligible, about $600$ at $x=57$ and $3000$ at $x=266$ for $c=5$, $c_0=-5$, in addition to the $400$ or so of the other parts.  For $c_0=0$ the general machinery reproduces the values of the specialized code to $10^{-13}$.

\paragraph{The fixed start}  The paths and partitions are those of the local null value, with the constant $\frac12$ of the parabola replaced by $\frac12(1+\gamma_0^2)$ in case \Fk, and with the exponent that places the nodes, $-tv^2-s_dX$, augmented by $\log\cosh\{b_2(2\tilde X)^{1/2}\}$, the growth of (32) and (33) away from the edge, whose derivative also enters the descent directions; the factor carried by the nodes is $\varphi_t$ divided by that growth, which is bounded.  Laid out from the first two terms alone, the paths miss the mass of $\varphi_t$ beyond the breakpoint, and the result depends on $\tau$ at the level of $10^{-11}$.  The half-line beyond the last breakpoint is the ray in the direction of the Gaussian's decay, with the pieces and node counts from the full exponent, or, when the exponential grows along it by more than $e^2$, the segment to the saddle of the quadratic exponent and the descent line from it.  The integral $G$ of (33) is a Gauss--Legendre rule in $T$ after $l=|2Y|^{1/2}\sinh T$, with $40$ nodes up to where the Gaussian, against the growth of the hyperbolic cosines, is below $e^{-41}$; $40$ nodes agree with $384$ to $10^{-14}$.  The Gaussian times the growth $e^{b_1l}$ peaks at $l^*=b_1/\Rea a$ with the width $(\Rea a)^{-1/2}$, and when that peak is separated from the scale $|2Y|^{1/2}$ of the edge, at large noncentrality, the rule in $T$ stops at $l^*-8(\Rea a)^{-1/2}$ and a Gauss--Legendre rule in $l$, with the node count from the change of the phase, covers the peak; the terms are formed relative to the largest, $G=Se^{E}$, so that nothing overflows, and $\cosh z$ as $e^z(1+e^{-2z})/2$ with $\Rea z\ge0$.  The rule runs along the ray $l=e^{-\iu\arg(a)/2}s$, on which the Gaussian of $G$ is real, so that its branch cut in $Y$ lies where $aY$ is negative real, as the cut of $K_0$ in (23) does, and the paths, which avoid that cut, serve.  The same mass that moves the peak of $G$ moves the mass of $\varphi_t$ beyond the breakpoint: on the ray in the direction of the Gaussian's decay the integrand there oscillates with a magnitude of up to $e^{\lambda}$ around a value of order one, so when the full exponent rises along the real axis by more than $2$ the ray is replaced by the path through the saddle of the full exponent (Newton from the maximum on the real axis, itself refined by Newton), continued by straight pieces in the direction of steepest descent, each renewed where the descent along it ends; the saddle is used only when it is that maximum slightly displaced, within one unit in height and with an imaginary part at most half its real part, since a saddle elsewhere is not one the real-axis path crosses: otherwise the ray is kept for a rise below $10$, where its cancellation is harmless, and the real-axis maximum is the pivot of the path beyond.  With $c>0$ the initial-value term $e^{\gamma_0^2Q(0,t)}$ has an essential singularity at $t_*$, and the line stays at $\tau\le t_*/2$, where $\tau=t_*/4$ gives the same values to $10^{-15}$.  In the C implementation the rule for $G$ is written in real arithmetic on arrays, with exponentials and sines from branch-free polynomial routines that the compiler vectorizes; a value of the distribution function then takes $7$ to $80$ milliseconds on ten threads of a laptop at $\lambda\le12$, against $0.25$ seconds by the conditional integral, and up to $0.3$ seconds at $\lambda=200$.

\paragraph{The moments}  The derivatives in (34) are Cauchy integrals,
\[
m_p(s)=\frac{p!}{Nr^p}\sum_{k=0}^{N-1}M(re^{\iu\theta_k},-s-c_0re^{\iu\theta_k})\,e^{-\iu p\theta_k},\qquad\theta_k=\frac{2\pi k}{N},
\]
with $N=64$ and $r=0.4\min_k|s_k(-s)|$, the roots of $D(\cdot,-s)$ being real; a circle on which $\log D$ does not return to its starting value encloses a zero of $D$, and its radius is halved.  The $s$-integral is the trapezoidal rule in $\log s$: the integrand is smooth in $\log s$ and decays at both ends, like $s^{p/2}$ at the left and like $\exp(-s^{1/2})$ at the right, so the rule converges geometrically.

\subsection{The conditional integral representation}
\label{app:double}

\paragraph{The conditional integral representation}  For each $v$ the inner inversion runs through a saddle point.  With a zero start the integrand is the product (\ref{eq:dkprod}) or $P(t)\varphi_t(a(v))e^{-tv}$ in the Girsanov forms, analytic beyond $t_*$; its saddle point $\hat s$ on the real axis is found by a safeguarded secant iteration, and the integral is taken along two paths through $\hat s$, a parabola that turns into a wedge and a vertical line in the $\vartheta$-plane, a parabola in $t$ opening to the right, and accepted when they agree; in case \Fk\ a third route writes $\erf=1-\erfc$ and inverts the density of $V$ and the remainder separately, each through its own saddle point.  With an initial observation the integrand is $M_V$ times $\cP$ of (\ref{eq:Pnc}), or times the tilted density, and the path is the parabola $t(\theta)=s^*+\mu\theta^2-2\iu\mu\theta$ of the Talbot method \citep{WeidemanTrefethen2007} through the saddle point $s^*$ of $\log M_V(t)-tv$, which opens to the right around the real singularities of $M_V$, with $\mu$ at least $4.5$, so that $\Rea(-z)^{1/2}\ge3$ along it, the branch of $\log D$ continued from $s^*$, and the sign conditions of Theorem~\ref{th:double} checked at every node.  Each inner integral uses the trapezoidal rule with $96$ nodes after a sinh substitution, with the nested rule of twice the step as error estimate, and the node count is doubled where the estimate requires it.  The outer integral over $v$ uses Gauss--Legendre rules in $\log v$ over $[e^{-8}\E V,e^6\E V]$ in parts of length at most two, with breakpoints at the zeros of $a(v)$, where a cosine map removes the singularities of the integrand, graded toward the sharp step of the integrand in the far upper tail for $c>c_0$, and bisected where an error estimate from the Legendre coefficients of the integrand exceeds $10^{-12}$; in (\ref{eq:Flim}) the boundary term becomes $\PP\{R\le a(v_0)\}$ at the lower end $v_0$.  The programs return an error estimate with each value and warn when it exceeds $10^{-10}$.


\paragraph{Zero start: the inversions}  For each pair $(v,a(v))$ let $\ell(t)$ be the logarithm of the integrand, (\ref{eq:dkprod}) in case \Fk\ and $\log\{P(t)\varphi_t(a(v))\}-tv$ in cases \Fu\ and \Ft, in the Girsanov forms of Section~3.1, with the exponent $-a\eta-tv$ of (\ref{eq:dkprod}) written as $-(a-cv)\eta-cv\epsilon+\delta^2v/2$, $\delta=\vartheta-c$ and $\epsilon=\eta-\delta$, which avoids the cancellation between $a\eta$ and $tv$ at large $v$.  We find its saddle point $\hat s$ on the real axis, left of the first point where one of $\cosh\vartheta\,|\bSig|$, $\alpha_0$ and $\alpha_1^2+4\alpha_0\alpha_2$ leaves $(0,\infty)$ (found by bisection), by bracketing and the Illinois variant of regula falsi on $\ell'$, with finite-difference steps scaled by the width $\ell''(\hat s)^{-1/2}$, and integrate along two paths through $\hat s$.  The first is
\[
t(\tau)=\hat s+\mu\frac{\tau^2}{(1+\tau^2/16)^{1/2}}-2\iu\mu\tau,\qquad\mu=\frac1{2\ell''(\hat s)^{1/2}},
\]
a parabola near $\hat s$ that turns into the wedge $\Rea t-\hat s=2|\Ima t|$ far out.  The second is the vertical line $\vartheta=\vartheta_A+\iu y$ in the $\vartheta$-plane, a parabola $t=(c^2-\vartheta^2)/2$ opening to the right, along which $|e^{-tv}|$ decays like $e^{-y^2v/2}$; we take $\vartheta_A=\max\{\hat\vartheta,\frac12\operatorname{asinh}(a/(6v)^{1/2}),1\}$, with $\hat\vartheta$ the value at $\hat s$, so that the oscillating part $a\vartheta(\coth\vartheta-1)$ of $a\eta$ stays small against $y^2v/2$.  In case \Fk\ a third route writes $\erf=1-\erfc$ and inverts the density of $V$ and the remainder separately, each along a path through its own saddle point left of $t_*$, with $\erfc$ from the Faddeeva function \citep{Weideman1994}; it is the better one when $\cP$ is close to one.  Each path uses the trapezoidal rule with $96$ nodes after a sinh substitution, with the nested rule of twice the step as error estimate, and its node count is doubled up to three times until that estimate is below $10^{-13}\max(1,|I|)$, $I$ the value.  The nested rule alone is not a safe estimate: an oscillating tail that neither the rule nor its half resolves can alias into the same value on both, and in case \Fk\ at $c=5$, $c_0=-5$, $v=0.086$ and $a=35$ the parabola with $96$ and with $192$ nodes agrees with itself to $10^{-15}$ while both are off by $10^{-7}$, which the vertical line, correct to $10^{-15}$ with $96$ nodes, exposes.  A value is therefore accepted when two paths agree to the tolerance, the one with the smaller estimate being kept and the discrepancy entering the error, or when the other paths are consistent with the one of smallest estimate, their discrepancies within ten times the sums of the estimates (a path that converges slowly, or not at all, is no evidence against a resolved one); otherwise the inconsistent paths and the best one are recomputed with the node count doubled, up to three more times, and if they still disagree the best value is returned with the smallest discrepancy as its error, which the outer rule reports.  For $G$ in (\ref{eq:Glim}) the crossing point is kept at least one saddle-point width from the pole at $t=0$, and the residue is added when it lies to the left.  The distribution function in case \Fk\ comes from part (i) and its density from (\ref{eq:flim}), whose integrand $e^{-c/2}(\vartheta/\sinh\vartheta)^{1/2}e^{-a\eta}(\pi a)^{-1/2}$ is the derivative of (\ref{eq:dkprod}) in $a$; in cases \Fu\ and \Ft\ both come from part (ii).

\paragraph{The outer integral}  The outer integral over $v$ uses Gauss--Legendre rules in $\log v$ on the interval from $v_0$ to $e^6\E_{\Fk}V$, where
\[
\E_{\Fk}V=\frac{e^{2c}-1-2c}{4c^2},
\]
the mean of $\int_0^1J(r)^2\dr{r}$, bounds $\E V$ in all three cases, and $v_0=3\times10^{-3}$ in cases \Fu\ and \Ft\ and the smaller of $3\times10^{-3}$ and $e^{-8}\E V$ in case \Fk; lowering $v_0$ to $10^{-5}$ changes the results by less than $10^{-10}$ for $c$ down to $-50$.  The parts are of length at most two, with $32$ nodes in case \Fk\ and $64$ in cases \Fu\ and \Ft, and with breakpoints at the zeros of $a(v)$, where the variable is mapped by $z\mapsto\frac12(1-\cos\pi z)$ (twice in cases \Fu\ and \Ft, for the logarithmic singularity of $\varphi_t$ at $r=0$), so that the distance to the breakpoint behaves like the fourth power of the variable.  For $c_0>0$ and $x$ near $-(2c_0)^{1/2}$ the two zeros of $a(v)$ merge, and the parts are graded geometrically toward them; at $x=-(2c_0)^{1/2}$ the density in case \Fk\ is infinite and is returned as such.  Two further devices handle the far upper tail for $c>c_0$.  There $R$ is close to $(c-c_0)V$, and the outer integrand has a sharp step in $\log v$, and the density a narrow bump, near the point where $a(v)$ crosses the conditional mean of $R$ given $V=v$.  We compute that point and the width of the step from the tilt at the saddle point of $\log M_V(s)-sv$: the width is the conditional standard deviation of $R$, from the double saddle-point formula $K_{uu}-K_{us}^2/K_{ss}$ for the conditional variance, divided by the slope of $a(v)$ minus the conditional mean, and when it is below $0.1$ the parts are graded geometrically toward the point, starting from twice the width.  In addition, a part is bisected, up to twelve times, when an error estimate from the Legendre coefficients of its integrand exceeds $10^{-12}$ and ten times the noise level of the inner integrals; the estimate is the larger of the last two coefficients, extrapolated by their decay from the middle ones.  With the fixed rule alone, the distribution function in case \Fk\ at $c=5$ is off by $3\times10^{-3}$ at $x=126$, with no warning, and the density is lost beyond $x\approx320$.  In (\ref{eq:Flim}) the boundary term becomes $\PP\{R\le a(v_0)\}$ at the lower end $v_0$, with the weights at $t=0$.  Pairs $(v,a(v))$ at which the saddle-point approximation of the inner integral, times the outer weight, is below $e^{-35}$ are skipped.  The outer estimates, plus the outer rule applied to the error estimates of the inner integrals, are returned as an error estimate, and values whose estimate exceeds $10^{-10}$ are flagged.

\paragraph{With an initial observation}  The inner integrals are one-dimensional inversions of $M_{\gamma_0}(0,t)\cP$ (fixed start, through (\ref{eq:Pnc})) or of $M_V$ times the tilted density (stationary start).  On the line itself the integrand decays only like $\exp(-|y|^{1/2}/2)$ in $y=\Ima t$, while $e^{-tv}$ makes it oscillate with period $2\pi/v$, so a direct rule needs many nodes, and we use the parabola $t(\theta)=s^*+\mu\theta^2-2\iu\mu\theta$, $\theta\in(-\infty,\infty)$, through the saddle point $s^*(v)$ of $\log M_V(s)-sv$ on the real axis (safeguarded Newton iteration, which keeps a bracket because $\log M_V$ is convex on its real domain), which opens to the right around the real singularities of $M_V$.  Along it $|e^{-tv}|=e^{-s^*v}e^{-\mu v\theta^2}$.  We take $\mu=\max\{1/(2\ell_2^{1/2}),4.5\}$, $\ell_2$ the second derivative of $\log M_V$ at $s^*$; the first choice turns the saddle-point Gaussian into $e^{-\theta^2/2}$, and the bound keeps $\Rea(-z)^{1/2}\ge3$ along the path, so that $e^{-2\vartheta}$ is small there.  The sign conditions of Theorem~\ref{th:double} do not hold on the whole parabola: in case \Fk, with $w=z^{1/2}=u+\iu v'$, $\kappa_1=1/\{2(w\cot w-c)\}$ and $\Rea\kappa_1>0$ if and only if $u\sin2u+v'\sinh2v'>c(\cosh2v'-\cos2u)$, which fails near $u=3\pi/4+k\pi$ once $u$ exceeds about $(2\mu)^{1/2}\sinh\{2(2\mu)^{1/2}\}$, at $|\theta|\approx202$ for $c=0$ and $\mu=4.5$.  We therefore use the parabola only for $|\theta|\le\Theta$, where $e^{-\mu v\Theta^2}<10^{-18}$, and close the path by the vertical rays $\Rea t=s^*+\mu\Theta^2$, whose contribution we neglect; the cutoff, like the range of the outer integral, is a practical choice, tested numerically: doubling the exponent at the cutoff and widening the range of $v$ to $[e^{-12},e^9]\E V$ changes the distribution function by less than $10^{-11}$.  In all our computations $\Theta\le9$, far inside the region where the conditions hold, and the routine checks them at every node it uses and on a grid of intermediate parabolas on request.  Along the parabola $\log M_V=-\frac12\log D(0,z)$ is continued from $s^*$, where $D>0$: we write $\log D=c+\vartheta+\log q$ with $\vartheta=(-z)^{1/2}$, continuous on the path because $\Ima(-z)>0$ there, and $q=De^{-c-\vartheta}$, which varies slowly and whose logarithm is unwrapped on a grid refined until the phase changes by less than $\pi/4$ between neighbors.  The real part of the integrand is even in $\theta$, and the trapezoidal rule is applied to $\theta\ge0$ after the substitution $\theta=\theta_w\sinh\tau$, with $\theta_w=1/(2\mu\ell_2^{1/2})$ the width of the saddle-point Gaussian, with $96$ nodes up to the point where $e^{-\mu v\theta^2}<10^{-18}$, doubled up to twice where the nested-rule estimate exceeds $10^{-9}$ of the largest column of the piece; with a fixed start in cases \Fu\ and \Ft\ the initial-value term has an essential singularity at the end of the domain of $M_{\gamma_0}(0,\cdot)$, toward which $s^*$ moves as $v$ grows, and the crossing point is kept at a relative distance of $2\%$ from the end.  Only an approximate $s^*$ is needed: any real point of the domain gives the same integral, and $s^*$ and $\ell_2$ only place and scale the path.  The boundary term $\PP(R\le\frac12)$ of (\ref{eq:Flim}) has real weights and is computed once for each $x$ from (\ref{eq:Pkappa}) or (\ref{eq:Pnc}), by $96$-point Gauss--Legendre after the substitution $l=\sqrt{2a}\sinh\tau$ when $a>0$ and $l=(u^2-2a)^{1/2}$ when $a<0$, as is the inner integral over $l$ of (\ref{eq:Pnc}) at every node.  The outer integral is the rule above in $\log v$ over $[e^{-8},e^6]\E V$ with $64$ nodes per part.

\paragraph{Accuracy of the conditional integral representation}  Doubling the outer and the inner node counts changes the distribution function and the density by at most $5\times10^{-11}$ with a zero start ($c\in[-30,10]$, $c_0\in[-40,10]$) and by $1.5\times10^{-9}$ at $c=c_0=20$.  In the far upper tail at $c=5$, where the outer integrand has the sharp step described above, the density integrated over $[-80,500]$ by a $400$-point Gauss--Legendre rule has mass $1+2\times10^{-9}$ in case \Fk\ and $1+2\times10^{-10}$ in case \Fu, and reproduces the means of Proposition~3 to $10^{-8}$, the remaining differences being the truncation and the rule of that integral.  With an initial observation, in case \Fk\ at $x=0$, where the exact value is a noncentral $\chi^2_1$ probability, the parabola route is accurate to $10^{-13}$ for $c\le2$; in cases \Fu\ and \Ft\ with a fixed start it agrees with the limit of the Richardson-extrapolated finite-sample initial-value terms of \cite{BaoBrodaKan2026} to $10^{-7}$, the accuracy of that extrapolation, and the sign conditions held at every node used for $c\le2$.

\subsection{Cost and accuracy}
\label{sec:cost}

\paragraph{Cost}  Both representations need about $10^5$ evaluations of the integrand per value, and their costs differ by the contour alone.  All timings are on one M4 laptop.  In the C implementation, on its ten threads, a value of the distribution function or the density by the line integral takes $2$ to $6$ milliseconds at $c_0=0$ in all six cases, $6$ to $9$ milliseconds with $c_0\ne0$, rising to $34$ in the upper tail at $c=5$, $c_0=-5$, where the saddle-factor part needs three outer nodes per radian of its phase, and $11$ to $110$ milliseconds with a fixed start; the conditional integral takes $4$ to $12$ milliseconds with a zero start for $c\le2$, $27$ to $39$ at $c=5$, $15$ to $16$ with a stationary start and $0.23$ to $0.26$ seconds with a fixed start (Tables~\ref{tab:acc} and~\ref{tab:accfixed}); on one thread both are three to four times slower.  The MATLAB programs, on one thread, take $0.03$ to $0.09$ seconds per value by the line integral at $c_0=0$ (medians over $x\in[-3,1]$), $0.15$ to $0.25$ seconds with $c_0\ne0$, rising to $2.6$ seconds at $x=342$ in the upper tail at $c=5$, $c_0=-5$, and by the conditional integral $0.05$ to $0.25$ seconds with a zero start ($0.8$ in case \Ft\ at $c=2$, $c_0=1$), $0.2$ to $0.5$ with a stationary start and $0.06$ to $0.12$ (case \Fk) and $1.4$ to $1.5$ seconds (cases \Fu\ and \Ft, where it is a triple integral) with a fixed start; a slower machine scales all of these alike.  The contour of the line integral is the same for every $x$, so that $P$ and the roots $s_1$, $s_2$ along it depend on $c$ alone and are shared between values of $x$; the Julia implementation exploits this and takes $1$ to $2$ milliseconds per value at $c_0=0$.  A quantile, by Newton's method on the distribution function with the density as derivative, takes a few values.

\paragraph{Accuracy}  Doubling every node count of the line integral changes the distribution function by at most $4\times10^{-13}$ and the density, relative to its value, by $10^{-11}$, over $x\in[-4,3]$ for the three cases at $c=0$ and $c=-5$; by at most $2\times10^{-12}$ at $c=-10$ and $-20$ around the center of the distribution; by at most $3\times10^{-12}$ for $c=1,2,3,5$ over $x\in[-4,80]$, where the density integrates to one to $10^{-9}$ and moving the line to $\tau=t_*/4$ changes nothing beyond $3\times10^{-12}$; and by at most $6\times10^{-13}$ in the stationary cases at $c=-1$, $-5$ and $-20$.  The trapezoidal rule in $w$ with the step $0.15$ applied to the undivided integrand agrees to the same accuracy.  Tables~\ref{tab:acc} and~\ref{tab:accfixed} set the two representations against each other, both in the C implementation, over the six cases, both sides of the unit root, the local null values of the checks of \cite{BaoBrodaKan2026}, and the stationary and fixed starts.  The distribution functions agree to $2\times10^{-11}$ at worst, except in the lower tails at $c=5$, where the line integral's absolute error reaches $3\times10^{-9}$ (case \Ft\ with $c_0=2$, at $x=-14$, where $\cF=3\times10^{-9}$) and $2\times10^{-10}$ (case \Fu\ with $c_0=-5$), values that the programs take from the conditional integral in any case (the tail rule below).  The densities agree to $3\times10^{-10}$ relative with a zero or stationary start and to $3\times10^{-9}$ with a fixed start, the largest differences at the points $2.5$ standard deviations out; this is the accuracy of the conditional integral itself.  The tables leave out the points at which the conditional integral reports that its error estimate is not met; the line integral is unchanged under node doubling at all of them.  They are case \Ft\ at $c=5$, $c_0=0$ and $x=33$, and, with a fixed start, case \Fk\ at $x=-3.8$ with $\gamma_0=3$, $c=-5$, at $x=-3.8$ and $-3.2$ with $\gamma_0=1$, $c=-20$, and case \Fu\ at $x=-4.3$ with $\gamma_0=1$, $c=-20$.  With a fixed start, at $18$ points over the three cases, $c\in[-5,2]$, $c_0=0$ and $c_0\ne0$, the two agree to $1.7\times10^{-13}$ in the distribution function and to $4\times10^{-11}$ in the density, again the accuracy of the conditional integral, and to $7\times10^{-14}$ at $c=-20$ and $-5$ with $\gamma_0$ up to $3$, where $\lambda$ reaches $200$; moving the contour of the line integral between $\tau=0.5$ and $2.5$ changes its values by $5\times10^{-15}$.  Beyond $c=2$, where the conditional integral no longer serves, the line integral with a fixed start is unchanged under node doubling to $10^{-16}$ and under $\tau\to\tau/2$ to $10^{-15}$ at $c=3$, $5$ and $10$, and in case \Fk\ it reproduces the closed form of $\PP(\tau\le0)$ at $c=5$ to $3\times10^{-16}$.  The error of the line integral is absolute rather than relative: the terms it sums are of order one, and it carries an error of $10^{-14}$ to $10^{-21}$ whatever the size of the value, which finer rules do not reduce ($2\times10^{-14}$ in the lower tails of cases \Ft\ and \St\ with $c_0\ne0$, and up to $3\times10^{-9}$ in the lower tails at $c=5$, where the saddle factor is large).  The programs therefore recompute by the conditional integral every value of the line integral below $10^{-3}$, a density or a distribution function in the lower tail, so that the relative accuracy holds in the tails as well.  Far out, below the error of either method, a computed density can be a tiny negative number; the programs return it as zero.

\begin{table}[tbp]
\caption{The two representations against each other}
\label{tab:acc}
\begin{center}
\footnotesize
\setlength{\tabcolsep}{4pt}
\begin{tabular}{rrcccccccc}
\toprule
 & & \multicolumn{3}{c}{$\max|\cF_{\rm line}-\cF_{\rm cond}|$} & \multicolumn{3}{c}{$\max|f_{\rm line}/f_{\rm cond}-1|$} & \multicolumn{2}{c}{time (s)}\\
\cmidrule(lr){3-5}\cmidrule(lr){6-8}\cmidrule(lr){9-10}
$c$ & $c_0$ & \Fk & \Fu & \Ft & \Fk & \Fu & \Ft & line & conditional\\
\midrule
$0$ & $0$ & $1\cdot10^{-14}$ & $6\cdot10^{-13}$ & $1\cdot10^{-12}$ & $1\cdot10^{-12}$ & $2\cdot10^{-12}$ & $4\cdot10^{-12}$ & 0.003 & 0.012\\
$-5$ & $0$ & $1\cdot10^{-14}$ & $3\cdot10^{-12}$ & $3\cdot10^{-12}$ & $1\cdot10^{-12}$ & $5\cdot10^{-12}$ & $2\cdot10^{-11}$ & 0.003 & 0.010\\
$-20$ & $0$ & $9\cdot10^{-14}$ & $2\cdot10^{-12}$ & $3\cdot10^{-12}$ & $1\cdot10^{-11}$ & $2\cdot10^{-10}$ & $1\cdot10^{-10}$ & 0.003 & 0.012\\
$2$ & $0$ & $6\cdot10^{-16}$ & $9\cdot10^{-13}$ & $8\cdot10^{-13}$ & $1\cdot10^{-12}$ & $2\cdot10^{-12}$ & $2\cdot10^{-12}$ & 0.002 & 0.005\\
$5$ & $0$ & $1\cdot10^{-15}$ & $5\cdot10^{-14}$ & $2\cdot10^{-11}$ & $1\cdot10^{-14}$ & $5\cdot10^{-14}$ & $4\cdot10^{-12}$ & 0.006 & 0.034\\
$0$ & $-5$ & $7\cdot10^{-16}$ & $7\cdot10^{-13}$ & $1\cdot10^{-12}$ & $6\cdot10^{-12}$ & $3\cdot10^{-12}$ & $4\cdot10^{-12}$ & 0.008 & 0.006\\
$2$ & $1$ & $7\cdot10^{-16}$ & $8\cdot10^{-13}$ & $5\cdot10^{-13}$ & $4\cdot10^{-12}$ & $2\cdot10^{-12}$ & $1\cdot10^{-12}$ & 0.007 & 0.006\\
$-2.5$ & $-5$ & $1\cdot10^{-15}$ & $2\cdot10^{-12}$ & $1\cdot10^{-12}$ & $6\cdot10^{-13}$ & $4\cdot10^{-12}$ & $3\cdot10^{-12}$ & 0.008 & 0.004\\
$5$ & $2$ & $3\cdot10^{-15}$ & $7\cdot10^{-14}$ & $3\cdot10^{-9}$ & $1\cdot10^{-10}$ & $1\cdot10^{-11}$ & $1\cdot10^{-12}$ & 0.006 & 0.027\\
$5$ & $-5$ & $5\cdot10^{-13}$ & $2\cdot10^{-10}$ & $9\cdot10^{-14}$ & $5\cdot10^{-11}$ & $2\cdot10^{-12}$ & $1\cdot10^{-12}$ & 0.034 & 0.039\\
$-10$ & $-20$ & $3\cdot10^{-15}$ & $3\cdot10^{-12}$ & $1\cdot10^{-12}$ & $6\cdot10^{-12}$ & $5\cdot10^{-12}$ & $2\cdot10^{-11}$ & 0.006 & 0.008\\
$2$ & $2$ & $1\cdot10^{-14}$ & $8\cdot10^{-13}$ & $5\cdot10^{-13}$ & $2\cdot10^{-13}$ & $1\cdot10^{-12}$ & $2\cdot10^{-12}$ & 0.009 & 0.006\\
$-2$ & $-2$ & $9\cdot10^{-16}$ & $1\cdot10^{-12}$ & $1\cdot10^{-12}$ & $1\cdot10^{-12}$ & $3\cdot10^{-12}$ & $2\cdot10^{-12}$ & 0.008 & 0.012\\
$-5$ & $3$ & $7\cdot10^{-16}$ & $3\cdot10^{-12}$ & $2\cdot10^{-12}$ & $8\cdot10^{-13}$ & $6\cdot10^{-12}$ & $4\cdot10^{-12}$ & 0.009 & 0.011\\
\midrule
 & & \Sk & \Su & \St & \Sk & \Su & \St & &\\
\cmidrule(lr){3-5}\cmidrule(lr){6-8}
$-5$ & $0$ & $2\cdot10^{-12}$ & $2\cdot10^{-12}$ & $2\cdot10^{-12}$ & $3\cdot10^{-12}$ & $4\cdot10^{-12}$ & $7\cdot10^{-11}$ & 0.003 & 0.016\\
$-20$ & $0$ & $7\cdot10^{-12}$ & $7\cdot10^{-12}$ & $7\cdot10^{-12}$ & $2\cdot10^{-11}$ & $4\cdot10^{-11}$ & $2\cdot10^{-10}$ & 0.004 & 0.015\\
\bottomrule
\end{tabular}
\end{center}
{\small Zero start (first block) and stationary start (second).  The largest differences between Theorem~2 and Theorem~\ref{th:double} over six or seven points per configuration, at the mean plus $-2.5$ to $2.5$ standard deviations; the relative density differences are over the points with $f>10^{-8}$.  Points at which the conditional integral reports that its error estimate is not met are left out (Section~\ref{sec:numerics}).  Time: the median per value of the distribution function by the C implementation on ten threads of an M4 laptop.}
\end{table}

\begin{table}[tbp]
\caption{As Table~\ref{tab:acc}, with a fixed start $\gamma_0$.}
\label{tab:accfixed}
\begin{center}
\footnotesize
\setlength{\tabcolsep}{4pt}
\begin{tabular}{rrrcccccccc}
\toprule
 & & & \multicolumn{3}{c}{$\max|\cF_{\rm line}-\cF_{\rm cond}|$} & \multicolumn{3}{c}{$\max|f_{\rm line}/f_{\rm cond}-1|$} & \multicolumn{2}{c}{time (s)}\\
\cmidrule(lr){4-6}\cmidrule(lr){7-9}\cmidrule(lr){10-11}
$c$ & $c_0$ & $\gamma_0$ & \Fk & \Fu & \Ft & \Fk & \Fu & \Ft & line & conditional\\
\midrule
$-5$ & $0$ & $1$ & $8\cdot10^{-16}$ & $2\cdot10^{-15}$ & $6\cdot10^{-14}$ & $3\cdot10^{-12}$ & $5\cdot10^{-12}$ & $4\cdot10^{-10}$ & 0.078 & 0.251\\
$-5$ & $0$ & $3$ & $2\cdot10^{-12}$ & $5\cdot10^{-15}$ & $8\cdot10^{-15}$ & $2\cdot10^{-11}$ & $6\cdot10^{-14}$ & $5\cdot10^{-12}$ & 0.087 & 0.247\\
$2$ & $0$ & $1$ & $2\cdot10^{-15}$ & $9\cdot10^{-16}$ & $2\cdot10^{-14}$ & $2\cdot10^{-12}$ & $8\cdot10^{-10}$ & $1\cdot10^{-9}$ & 0.011 & 0.261\\
$-20$ & $0$ & $1$ & $1\cdot10^{-13}$ & $2\cdot10^{-13}$ & $2\cdot10^{-14}$ & $1\cdot10^{-11}$ & $5\cdot10^{-12}$ & $5\cdot10^{-12}$ & 0.110 & 0.227\\
$-5$ & $-5$ & $1$ & $1\cdot10^{-15}$ & $2\cdot10^{-15}$ & $5\cdot10^{-14}$ & $7\cdot10^{-12}$ & $1\cdot10^{-11}$ & $3\cdot10^{-9}$ & 0.076 & 0.237\\
\bottomrule
\end{tabular}
\end{center}
\end{table}


\subsection{Checks against known values and the literature}
\label{sec:checks}

The closed form (7) of the paper was checked against the limit of the finite-sample characteristic function of \cite{BaoBrodaKan2026}, obtained by Richardson extrapolation from $n=2^{20}$ and $2^{21}$, to $2\times10^{-8}$ in all three cases, at real and complex arguments and for $c\in\{-10,-2.5,0,1.5\}$, which is an independent check of the determinant of Model~4, which \cite{NabeyaTanaka1990} derived by computer algebra.  The distribution function reproduces the probability $\PP(\hat\tau\le0)$, which is known in closed form, to $10^{-13}$ for $c$ between $-30$ and $20$, the series and the integrals of \cite{Abadir1995} to $10^{-15}$, the probability $\PP(\hat\tau_\mu\le0)$ computed by \cite{BaoKan2026} from the distribution of the coefficient statistic to its eight digits, and the response-surface critical values of \cite{MacKinnon1996,MacKinnon2010} to their stated accuracy; the two representations agree with each other to $10^{-11}$ over the six cases, both sides of the unit root and the local null values; the moments reproduce the values of \cite{Nabeya1999} and \cite{Larsson1997}; the local powers reproduce the simulations of \cite{Larsson1995} and \cite{Elliott1999} within their sampling error; the powers with a fixed start of Table~IV of the paper were checked against Theorem~\ref{th:double} with (\ref{eq:Pnc}), which agrees to the three decimals shown; in the GLS-detrended trend case the programs reproduce the closed form of $\PP(\tau\le0)$ at $c=0$, $\erf\{(0.4)^{-1/2}(1-\lambda)^{-1}\}=0.9943906$, to $10^{-12}$; and two million simulated paths of the limiting functionals agree with the distribution functions within their standard errors.  The details follow.

\paragraph{Known values}  In case \Fk\ with $c_0=0$ the limit at $x=0$ is known: $\PP(\tau\le0)=\PP\{J(1)^2\le1\}=\erf\{(2w)^{-1/2}\}$ with $w=\Var J(1)=(e^{2c}-1)/(2c)$, which at $c=0$ is $2\Phi(1)-1$; with a fixed start, $\PP(\tau\le0)=\PP\{J(1)^2\le1+\gamma_0^2\}$ with $J(1)\sim N(\gamma_0e^c,w)$.  Both representations reproduce these to $3\times10^{-13}$ for $c\in\{-30,-10,-2.5,0,1,2,5,10,20\}$ with a zero start, where at $c=20$ the probability is $1.04\times10^{-8}$, and to $2\times10^{-15}$ with $\gamma_0=1.5$ at $c=-5$ and $c=2$.  In case \Fu\ they give $\PP(\tau_\mu\le0)=0.957565589880$, against the $0.95756559$ that \cite{BaoKan2026} obtain from the distribution of the coefficient statistic, since $\{\hat\tau_\mu\le0\}=\{\phih\le1\}$.  When $c_0\ne0$ the probability at $x=0$ is $\PP(U-c_0V\le0)$, which has an independent one-dimensional representation, $-(2\pi\iu)^{-1}\int_{\sigma-\iu\infty}^{\sigma+\iu\infty}M(\zeta,-c_0\zeta)\zeta^{-1}\dr{\zeta}$ with $\sigma<0$, that decays fast along the line because of the factor in $V$; at $(c,c_0)=(0,-5)$, $(2,1)$, $(-10,-20)$, $(-2.5,-5)$, $(5,2)$, $(5,-5)$, $(10,5)$ and $(10,-5)$ the three computations agree to $1.2\times10^{-11}$ in all three cases, and at $(20,10)$, where the probabilities are about $2\times10^{-8}$, to $10^{-16}$.  The stationary-start distribution function equals the fixed-start one averaged over $\gamma_0\sim N(0,\sigma_0^2)$, an identity that neither computation uses: with 24-point Gauss--Hermite quadrature over $\gamma_0$, at $c=-5$ and $x=-2$, the two agree to $8\times10^{-13}$, $8\times10^{-14}$ and $1.2\times10^{-13}$ in cases \Sk, \Su\ and \St, and as $\gamma_0\to0$ the fixed-start distribution functions tend to those of the zero start, to $4\times10^{-13}$ at $\gamma_0=10^{-6}$.  The closed form (7) agrees with the Richardson limit of the finite-sample moment generating function of \cite{BaoBrodaKan2026} to $4\times10^{-9}$, the accuracy of that extrapolation, and its partial fractions (10) fit the limit as a function of $s$ with no room for a further linear or quadratic term; at real arguments the squared means $\mu_k$ are positive, and the prediction of (10) at a third value of $s$ from two others holds to $10^{-16}$.  The densities agree with fifth-order central differences of the distribution functions to $4\times10^{-13}$.  Finally, two million paths of the Ornstein--Uhlenbeck process, by Euler's scheme with $1000$ steps, give the limiting functionals (3) directly: at $c=0$ and $c=-5$ with a zero start, and at $c=-5$ and $-10$ with a stationary start, the simulated distribution functions at $x=-3,-2,-1,0$ agree with the exact ones to within $0.003$ at all $48$ points, the differences being the discretization error of the scheme, which at this number of paths exceeds the sampling error of $3\times10^{-4}$ (for $\hat\tau_\tau$ at $c=0$ and $x=-2$, with $4\times10^5$ paths, the difference is $-0.013$, $-0.007$ and $-0.003$ with $250$, $500$ and $1000$ steps, and within the sampling error of $0.0008$ with $2000$ and $4000$); with a fixed start $\gamma_0=1.5$ and with a stationary start, $2\times10^5$ simulated regressions at $n=1000$ and $\phi=1-5/n$ agree with the limits within about one standard error at $x\in\{-3,-2,-1,0\}$ in all six cases.


\bibliographystyle{econsoc}
\bibliography{LTR_BBK}

\end{document}
